\documentclass[12pt, aspectratio=169]{beamer}

\input{../header}
\usepackage{booktabs}
% \usepackage{tikz}
% \usetikzlibrary{patterns}
% \DeclareMathOperator{\E}{E}
% \DeclareMathOperator{\Var}{Var}
% \DeclareMathOperator{\V}{V}
% \DeclareMathOperator{\C}{C}


\title{Common Sampling Distributions}
\author{Md. Aminul Islam Shazid}
\date{}


\begin{document}
    {
		\setbeamertemplate{footline}{}    % NO FOOTLINE FOR THESE TWO FRAMES
		\addtocounter{framenumber}{-2}    % not counting the title page and the outline in frame numbers

		\begin{frame}
			\titlepage
		\end{frame}

		\begin{frame}{Outline}
            \vfill
			\tableofcontents[subsectionstyle=hide]
            \vfill
		\end{frame}
	}


	\section{Introduction}
    
    \begin{frame}{Sampling Distributions}
        \begin{itemize}
            \item The value of a statistic varies across different samples
            \item A \textbf{sampling distribution} is the theoretical probability distribution of a statistic computed from repeated random samples of a population
            \item Common statistics include the sample mean, sample variance, and ratios of variances
            \item Sampling distributions describe how these statistics vary from sample to sample
            \item They form the theoretical basis for statistical inference and hypothesis testing
        \end{itemize}
    \end{frame}


    \begin{frame}{Why Sampling Distributions Matter}
        \begin{itemize}
            \item Population parameters are usually unknown
            \item Instead, we compute statistics from samples
            \item Sampling distributions allow us to:
            \begin{itemize}
                \item quantify sampling variability,
                \item calculate probabilities for observed statistics,
                \item construct confidence intervals,
                \item perform hypothesis tests.
            \end{itemize}
        \end{itemize}
    \end{frame}


    \begin{frame}{Degrees of Freedom}
        \begin{itemize}
            \item \textbf{Degrees of freedom (DoF)} refer to the number of independent pieces of information available when estimating a parameter

            \item In many statistical calculations, some information is used to estimate parameters. This reduces the number of values that are free to vary

            \item Example: suppose we have a sample of size \(n\) and we compute the sample mean \(\bar{X}\)

            \item Once the mean is fixed, only \(n-1\) observations can vary freely. The final observation is determined by the requirement that the sample mean remains unchanged

            \item Therefore, the sample variance uses \(n-1\) degrees of freedom

            \item Many sampling distributions depend on degrees of freedom, including Chi-square distribution, Student's \(t\) distribution, \(F\) distribution
        \end{itemize}
    \end{frame}


    \begin{frame}{Chi-Square \((\chi^2)\) Distribution}
        \begin{itemize}
            \item If \(Z_1, Z_2, \dots, Z_k\) are independent standard normal variables, then
            \[
                \chi^2 = Z_1^2 + Z_2^2 + \dots + Z_k^2
            \]
            follows a \textbf{Chi-Square distribution} with \(k\) degrees of freedom
            \item The distribution depends on the parameter called \textbf{degrees of freedom}
            \item It is defined only for positive values
            \item The distribution is right-skewed, especially for small degrees of freedom
        \end{itemize}
    \end{frame}


    \begin{frame}{Use of Chi-Square Distribution}
        \begin{itemize}
            \item Testing hypotheses about population variance
            \item Goodness-of-fit tests
            \item Tests of independence in contingency tables
            \item Many statistical procedures rely on the Chi-square distribution of the sample variance
        \end{itemize}
    \end{frame}


    \begin{frame}{Student's \(t\) Distribution}
        \begin{itemize}
            \item Let \(\bar{X}\) be the sample mean of the random variable \(X\). Then \(\bar{X}\) is also a random variable because the sample mean varies across different samples
            \item The statistic
            \[
                t = \frac{\bar{X} - \mu}{s/\sqrt{n}}
            \]
            follows a \(t\)-distribution with \(n-1\) degrees of freedom, where:
            \begin{itemize}
                \item \(\mu\) is the population mean
                \item \(s\) is the sample variance
                \item \(n\) is the sample size
            \end{itemize}
            \item The distribution is symmetric around zero
        \end{itemize}
    \end{frame}


    \begin{frame}{Use of \(t\) Distribution}
        \begin{itemize}
            \item Testing hypotheses about a population mean when variance is unknown
            \item Constructing confidence intervals for the mean
        \end{itemize}
    \end{frame}


    \begin{frame}{\(F\) Distribution}
        \begin{itemize}
            \item The \textbf{\(F\) distribution} arises from the ratio of two independent Chi-square variables
            \item If
            \[
                F = \frac{(U_1 / d_1)}{(U_2 / d_2)}
            \]
            where \(U_1 \sim \chi^2(d_1)\) and \(U_2 \sim \chi^2(d_2)\), then \(F\) follows an \(F(d_1,d_2)\) distribution
            \item The distribution depends on two degrees of freedom parameters
            \item It takes only positive values and is right-skewed
        \end{itemize}
    \end{frame}


    \begin{frame}{Use of \(F\) Distribution}
        \begin{itemize}
            \item Comparing two population variances
            \item Analysis of variance (ANOVA)
            \item Regression model significance testing
        \end{itemize}
    \end{frame}


    % \begin{frame}{Summary of Sampling Distributions}
    %     \begin{center}
    %         \begin{tabular}{lcc}
    %             \toprule
    %             Distribution & Main Use & Key Parameter \\
    %             \midrule
    %             Chi-square & Variance related tests & Degrees of freedom \\
    %             \(t\) & Mean inference (small sample) & Degrees of freedom \\
    %             \(F\) & Variance comparison / ANOVA & Two degrees of freedom \\
    %             \bottomrule
    %         \end{tabular}
    %     \end{center}
    % \end{frame}

    \section*{Thank you.}

\end{document}